% Proposition: group-normalised advantages (GRPO-style) with an in-loop filter (draft v1)
% Numerical companion: code/grpo_champ_moyen_v1.py, results resultats/grpo-v1/

\begin{proposition}[Group-normalised advantages only change the clock]\label{prop:grpo}
Consider the tabular three-stratum model with survival probabilities $\beta_{\mathrm{h}} = 1$, $\beta_{\mathrm{v}}, \beta_{\mathrm{f}} \in [0,1]$. Each training step draws groups of $G$ independent samples from $\pi_\theta$; a cheating sample survives the monitor with probability $\beta_i$ and then earns $\Delta R$, otherwise it earns $-C$ and stays in the group (penalty variant); honest samples earn $0$. The advantage of a sample is $(r - \bar r_G)/s_G$, with $\bar r_G, s_G$ the group mean and (population) standard deviation, and $0$ when $s_G = 0$. Let $\bar r_i := \beta_i R_i - (1 - \beta_i) C$, $\mu := \sum_j \pi_j \bar r_j$ and $\sigma^2 := \Var_{\pi}(r) > 0$ (variance of the per-sample reward, interception randomness included).
\begin{enumerate}
  \item As $G \to \infty$, $\E[\text{advantage} \mid \text{sample in stratum } i] \to (\bar r_i - \mu)/\sigma$ for every $i$ with $\pi_i > 0$.
  \item In this limit, the natural-gradient flow without KL anchoring, $\dot\theta_i = (\bar r_i - \mu)/\sigma$, has the explicit solution $\pi_t \propto \pi_0 \exp\big(\tau(t)\, \bar r\big)$ with $\tau(t) = \int_0^t \sigma_u^{-1}\, du$: the trajectory is \emph{exactly} the unnormalised operator-B path, traversed with the clock $\tau$. Standardisation changes when a threshold is crossed, never which policy the dynamics select.
  \item With $\beta_{\mathrm{f}} = 1$, $\sigma \ge \sqrt{\pi_{\mathrm{v}}\beta_{\mathrm{v}}(1-\beta_{\mathrm{v}})}\,(\Delta R + C)$, so the per-step log-odds rates are bounded uniformly in $C$: $0 < \dot\ell_{\mathrm{f:h}} = \Delta R/\sigma \to 0$ and $\dot\ell_{\mathrm{f:v}} = (1-\beta_{\mathrm{v}})(\Delta R + C)/\sigma \le \sqrt{(1-\beta_{\mathrm{v}})/(\pi_{\mathrm{v}}\beta_{\mathrm{v}})}$ for $\beta_{\mathrm{v}} > 0$. In particular, under group normalisation a larger penalty does not accelerate the growth of stealth relative to visible cheating without bound, and it slows the growth of stealth relative to honesty per step; in the clock $\tau$, the comparison is exactly that of Theorem~\ref{thm:selection}.
  \item (Masking variant.) If intercepted samples are removed from the group, the limit is $\E[\text{advantage}\,\mathbf{1}\{\text{kept}\} \mid i] \to \beta_i (R_i - \mu_{\mathrm{acc}})/\sigma_{\mathrm{acc}}$, with $\mu_{\mathrm{acc}}, \sigma_{\mathrm{acc}}$ the mean and standard deviation of the reward over surviving samples: the discard policy gradient of Theorem~\ref{thm:dynamics-tabular} with the clock $\int \sigma_{\mathrm{acc}}^{-1}$.
\end{enumerate}
\end{proposition}

\begin{proof}
\emph{Part 1.} Fix a sample in stratum $i$ and write its reward $r$ and the other $G-1$ rewards $r_2, \dots, r_G$ (i.i.d.\ from the per-sample reward law under $\pi$, independent of $r$). By the strong law of large numbers, $\bar r_G \to \mu$ and $s_G \to \sigma > 0$ almost surely, so the advantage converges almost surely to $(r - \mu)/\sigma$. The standardised value of one point of a sample of size $G$ is bounded in absolute value by $\sqrt{G-1}$, which is not uniform in $G$; instead use that rewards take values in the finite set $\{-C, 0, \Delta R\}$, so $|r - \bar r_G| \le \Delta R + C$, and on the event $\{s_G \ge \sigma/2\}$ the advantage is bounded by $2(\Delta R + C)/\sigma$. The complementary event has probability $O(e^{-cG})$ by Hoeffding's inequality applied to the first two empirical moments, and there the advantage is bounded by $\sqrt{G-1}$; hence $\E|\text{advantage}|\mathbf{1}\{s_G < \sigma/2\} \le \sqrt{G-1}\,O(e^{-cG}) \to 0$, and dominated convergence on $\{s_G \ge \sigma/2\}$ gives $\E[\text{advantage} \mid i] \to \E[(r - \mu)/\sigma \mid i] = (\bar r_i - \mu)/\sigma$.

\emph{Part 2.} For a tabular softmax, the natural gradient of an advantage-weighted objective is the vector of per-stratum expected advantages, up to a common additive constant (the Fisher pseudo-inverse acts as division by $\pi_i$ on vectors orthogonal to $\mathbf{1}$ under $\pi$). Since $\mu$ and $\sigma_t$ are common to all coordinates, $\theta_i(t) - \theta_i(0) = \bar r_i\, \tau(t) - \int_0^t \mu_u \sigma_u^{-1} du$, and the second term is common to all $i$ and cancels in the softmax.

\emph{Part 3.} With $\beta_{\mathrm{f}} = \beta_{\mathrm{h}} = 1$, only visible cheats are intercepted, and the law of total variance gives $\sigma^2 \ge \E_\pi[\Var(r \mid \text{stratum})] = \pi_{\mathrm{v}}\beta_{\mathrm{v}}(1-\beta_{\mathrm{v}})(\Delta R + C)^2$. Then $\dot\ell_{\mathrm{f:h}} = (\bar r_{\mathrm{f}} - \bar r_{\mathrm{h}})/\sigma = \Delta R/\sigma$ and $\dot\ell_{\mathrm{f:v}} = (\bar r_{\mathrm{f}} - \bar r_{\mathrm{v}})/\sigma = (1-\beta_{\mathrm{v}})(\Delta R + C)/\sigma$, which give the bounds; $\sigma \to \infty$ as $C \to \infty$ when $\pi_{\mathrm{v}}\beta_{\mathrm{v}}(1-\beta_{\mathrm{v}}) > 0$.

\emph{Part 4.} Identical to Part 1 with the group restricted to survivors; a sample of stratum $i$ is kept with probability $\beta_i$, independently of the rest, and contributes $0$ when masked.
\end{proof}

\paragraph{Numerical check.} \texttt{grpo\_champ\_moyen\_v1.py} estimates the per-stratum expected advantage by Monte Carlo for $\pi = (0.80, 0.15, 0.05)$, $\beta_{\mathrm{v}} = 0.1$, $\beta_{\mathrm{f}} \in \{1, 0.75\}$, $C \in \{0, 1, 4\}$: the maximal deviation from the mean-field value decreases from $0.56$--$2.2$ at $G = 4$ to $0.008$--$0.025$ at $G = 256$. With $\beta_{\mathrm{f}} = 1$, raising $C$ from $0$ to $10$ lowers the per-step rate $\dot\ell_{\mathrm{f:h}}$ from $4.06$ to $0.29$ and $\dot\ell_{\mathrm{f:v}}$ from $3.65$ to $2.87$; with $\beta_{\mathrm{f}} = 0.75$, $\dot\ell_{\mathrm{f:h}}$ changes sign at $C = 3$, the operator-B barrier $\beta_{\mathrm{f}}\Delta R = (1 - \beta_{\mathrm{f}})C$.
